%! End Behavior and Asymptotes — Unit 2, Lesson 2.5 — Homework (Answer Key)
\documentclass[10pt]{article}
\usepackage{atda-article}
\usepackage{atda-key}   % pulls in -boxes; do NOT also load -boxes
\newcommand{\TallMath}[1]{$\displaystyle #1\rule[-1.4em]{0pt}{3.2em}$}
\setlength{\workrowsep}{8pt}

\begin{document}

\pageheader{Unit 2, Lesson 2.5}{Homework --- Answer Key}
% NAMESTRIP: no \namedateperiod here — mirrors the blank. Teacher-only prose goes
% in the lesson plan as \begin{teachernote}[Homework], never in a key.

\vspace{0.15in}

\boxguard[24]
\begin{notesbox}{Practice --- The Ends of the Picture}
\small
Every asymptote you name gets written as an \emph{equation} of a line, never as a bare number.

\begin{enumerate}[leftmargin=1.6em, itemsep=6pt, topsep=4pt]

\item \textbf{End behavior from three graphs.} Fill the table. Each answer is one of: \emph{climbs
without bound}, \emph{falls without bound}, or \emph{levels off toward \underline{\hspace{0.7cm}}}.

\begin{center}
\begin{tabular}{@{}c@{\hspace{0.4cm}}c@{\hspace{0.4cm}}c@{}}
\begin{tikzpicture}
\begin{axis}[
    width=5.0cm, height=3.9cm, title={\scriptsize (a) $y=-x^2+4$},
    xmin=-3.4, xmax=3.4, ymin=-5.4, ymax=5.4,
    xtick={-3,-2,-1,0,1,2,3}, ytick={-4,-2,0,2,4}, minor y tick num=1,
    grid=both, grid style={linegray!45}, minor grid style={linegray!30},
    axis lines=left, tick label style={font=\tiny}]
  \addplot[royal, very thick, domain=-3:3, samples=90]{-x^2+4};
\end{axis}
\end{tikzpicture}
&
\begin{tikzpicture}
\begin{axis}[
    width=5.0cm, height=3.9cm, title={\scriptsize (b) $y=3^x$},
    xmin=-3.4, xmax=2.4, ymin=-1.4, ymax=10.4,
    xtick={-3,-2,-1,0,1,2}, ytick={0,2,4,6,8,10}, minor y tick num=1,
    grid=both, grid style={linegray!45}, minor grid style={linegray!30},
    axis lines=left, tick label style={font=\tiny}]
  \addplot[linegray, thick, dashed, domain=-3.3:2.3, samples=2]{0};
  \addplot[royal, very thick, domain=-3.3:2.1, samples=90]{3^x};
\end{axis}
\end{tikzpicture}
&
\begin{tikzpicture}
\begin{axis}[
    width=5.0cm, height=3.9cm, title={\scriptsize (c) $y=-2x+2$},
    xmin=-3.4, xmax=3.4, ymin=-4.4, ymax=8.4,
    xtick={-3,-2,-1,0,1,2,3}, ytick={-4,-2,0,2,4,6,8}, minor y tick num=1,
    grid=both, grid style={linegray!45}, minor grid style={linegray!30},
    axis lines=left, tick label style={font=\tiny}]
  \addplot[royal, very thick, domain=-3.3:3.3, samples=2]{-2*x+2};
\end{axis}
\end{tikzpicture}
\end{tabular}
\end{center}

\vspace{-2pt}
\renewcommand{\arraystretch}{1.4}
\begin{tabularx}{\linewidth}{@{}p{1.4cm} X X p{2.9cm}@{}}
\rowcolor{mist}
\textbf{Graph} & \textbf{Far left ($x \to -\infty$)} & \textbf{Far right ($x \to \infty$)} &
    \textbf{Horizontal}\newline\textbf{asymptote?} \\
\hline
(a) & \ans{falls without bound} & \ans{falls without bound} & \ans{none} \\
(b) & \ans{levels off toward $0$} & \ans{climbs without bound} & \ans{$y=0$} \\
(c) & \ans{climbs without bound} & \ans{falls without bound} & \ans{none} \\
\end{tabularx}

\item \textbf{The asymptote travels with the graph} (Lesson 2.2 spiral). The graph of $y=2^x$ has the
horizontal asymptote $y=0$. Give the horizontal asymptote of each function below as an equation.
\smallskip

\renewcommand{\arraystretch}{1.35}
\begin{tabularx}{\linewidth}{@{}X p{3.0cm} X p{3.0cm}@{}}
\rowcolor{mist}
\textbf{Function} & \textbf{Asymptote} & \textbf{Function} & \textbf{Asymptote} \\
\hline
$y=2^x+3$ & \ans{$y=3$} & $y=2^{x+4}$ & \ans{$y=0$} \\
$y=2^x-1$ & \ans{$y=-1$} & $y=-2^x$ & \ans{$y=0$} \\
\end{tabularx}

\item \textbf{Which one is it?} Label each description with a term from the bank.
\smallskip

{\footnotesize\textbf{Bank:} zero \quad$\cdot$\quad $y$-intercept \quad$\cdot$\quad end behavior
\quad$\cdot$\quad horizontal asymptote \quad$\cdot$\quad vertical asymptote}
\smallskip

\renewcommand{\arraystretch}{1.35}
\begin{tabularx}{\linewidth}{@{}X p{3.6cm}@{}}
\rowcolor{mist}
\textbf{The name for\dots} & \textbf{is\dots} \\
\hline
what the outputs do as the inputs run off both edges & \ans{end behavior} \\
a line the graph nears as the inputs run off to the right & \ans{horizontal asymptote} \\
an input at which the function has no value and the outputs blow up & \ans{vertical asymptote} \\
an input at which the output is exactly $0$ & \ans{zero} \\
the point where the graph crosses the vertical axis & \ans{$y$-intercept} \\
\end{tabularx}

\item \textbf{How many asymptotes?} For each function, say how many horizontal asymptotes and how
many vertical asymptotes its graph has, and give the reason in a few words. No graphing.
\begin{enumerate}[label=(\alph*), leftmargin=1.6em, itemsep=3pt, topsep=3pt]
  \item $y=x^2-5$
  \item $y=4^x$
  \item $y=3x+1$
  \item $y=\left(\tfrac12\right)^x$
\end{enumerate}

\ansline{(a) None and none. A parabola climbs without bound at both ends, and its domain is all reals.}\\
\ansline{(b) One horizontal, $y=0$ (it levels off going left); no vertical --- every input is allowed.}\\
\ansline{(c) None and none. A line falls without bound one way and climbs without bound the other.}\\
\ansline{(d) One horizontal, $y=0$ (it levels off going \emph{right} this time); no vertical.}\\

\item \textbf{A graph with a hole in its domain.} Use the graph of $h$ below.

\begin{center}
\begin{tikzpicture}
\begin{axis}[
    width=9.4cm, height=4.6cm,
    xlabel={$x$}, ylabel={$y$},
    xmin=-6, xmax=10, ymin=-8, ymax=8,
    xtick={-6,-4,-2,0,2,4,6,8,10}, ytick={-8,-6,-4,-2,0,2,4,6,8},
    minor x tick num=1, minor y tick num=1,
    grid=both, grid style={linegray!45}, minor grid style={linegray!30},
    axis lines=left,
    tick label style={font=\tiny}, label style={font=\scriptsize}]
  \addplot[linegray, thick, dashed, domain=-6:10, samples=2]{0};
  \addplot[linegray, thick, dashed] coordinates {(2,-8) (2,8)};
  \addplot[royal, very thick, domain=-6:1.25, samples=140]{6/(x-2)};
  \addplot[royal, very thick, domain=2.75:10, samples=140]{6/(x-2)};
  \addplot[keyred, only marks, mark=*, mark size=1.7pt] coordinates {(0,-3)};
\end{axis}
\end{tikzpicture}
\end{center}

Give the vertical asymptote as an equation, the horizontal asymptote as an equation, the domain of
$h$, and the end behavior at both ends. Then say why this graph has no $y$-intercept\dots{} or does
it? Check the picture before you answer.

\ansline{Vertical asymptote $x=2$; horizontal asymptote $y=0$. Domain: all real numbers except $2$.}\\
\ansline{As $x \to \infty$ the outputs level off toward $0$ from above; as $x \to -\infty$, toward $0$ from below.}\\
\ansline{It \emph{does} have a $y$-intercept: $(0,-3)$. Only $x=2$ is missing from the domain, and}\\
\ansline{$x=0$ is perfectly allowed. A vertical asymptote removes one input, not the whole axis.}\\

\end{enumerate}
\end{notesbox}

\vspace{0.10in}

\boxguard[30]
\begin{scenariobox}[In Context --- The Reservoir Refills]{royal}
\small
A drought left the town reservoir at $36\%$ of capacity. Each week the inflow refills half of
whatever is still missing. The graph shows $R(w)$, the percent of capacity $w$ weeks after the rains
began. The dashed line is a landmark.

\begin{center}
\begin{tikzpicture}
\begin{axis}[
    width=10.6cm, height=4.6cm,
    xlabel={$w$ (weeks after the rains began)}, ylabel={$R$ (percent of capacity)},
    xmin=0, xmax=6, ymin=0, ymax=110,
    xtick={0,1,2,3,4,5,6}, ytick={0,20,40,60,80,100}, minor y tick num=4,
    grid=both, grid style={linegray!45}, minor grid style={linegray!30},
    axis lines=left,
    tick label style={font=\tiny}, label style={font=\scriptsize}]
  \addplot[linegray, thick, dashed, domain=0:6, samples=2]{100};
  \addplot[royal, very thick, domain=0:6, samples=140]{100-64*(0.5)^x};
  \addplot[keyred, only marks, mark=*, mark size=1.7pt]
    coordinates {(0,36) (1,68) (2,84) (3,92) (4,96)};
\end{axis}
\end{tikzpicture}
\end{center}

\begin{enumerate}[leftmargin=1.6em, itemsep=6pt, topsep=2pt, start=6]

\item Read the graph and fill in the table. Then describe the end behavior as $w$ runs off to the
right.
\smallskip

\renewcommand{\arraystretch}{1.3}
\begin{tabular}{@{}|l|c|c|c|c|c|@{}}
\hline
$w$ (weeks) & $0$ & $1$ & $2$ & $3$ & $4$ \\ \hline
$R(w)$ (\% full) & \ans{$36$} & \ans{$68$} & \ans{$84$} & \ans{$92$} & \ans{$96$} \\
\hline
\end{tabular}

\ansline{As $w \to \infty$ the outputs level off toward $100$. The missing share halves each week:}\\
\ansline{$64, 32, 16, 8, 4, \dots$ --- always positive, so the graph keeps rising and never levels out flat.}\\

\item Give the horizontal asymptote as an equation and say what that line means about the reservoir.
Then name the interval on which $R$ is increasing (Lesson 2.4 spiral).

\ansline{Horizontal asymptote: $y=100$.}\\
\ansline{It is the reservoir's full capacity --- the level this model approaches and never quite reaches.}\\
\ansline{$R$ is increasing on $(0,\infty)$: it rises for every week, without ever changing direction.}\\

% Unguarded, item 8's stem and two of its three write-line slots sat at the foot of
% p2 with the third opening p3 -- an answer a student starts at the bottom of one
% page and finishes at the top of the next. \boxguard is inert inside a breakable
% tcolorbox, so \tcbbreak is the only lever. Mirrors homework.
\tcbbreak
\item Does $R$ have an absolute maximum? Does it have an absolute minimum? Give a value and a
location where one exists, and say why not where one does not.

\ansline{No absolute maximum. Every week is higher than the last, so no output is the highest one ---}\\
\ansline{and $100$ itself is never an output.}\\
\ansline{Absolute minimum $36\%$ at $w=0$ --- at an endpoint, the week the rains began.}\\

\item The water authority announces: \emph{``The reservoir will be completely full in 8 weeks.''}
Using this model, settle the claim. Then write a statement they could make instead that is both true
and useful to the town. (At $w=8$ the missing share is $0.25\%$.)

\ansline{False as stated. By this model the reservoir is never \emph{completely} full: it approaches}\\
\ansline{$y=100$ and never reaches it, so ``full in 8 weeks'' cannot happen in 8 weeks or ever.}\\
\ansline{True and useful: ``In 8 weeks the reservoir will be over $99.7\%$ of capacity.'' That is}\\
\ansline{$100-0.25$, and being a quarter of a percent short is not the same as being full.}\\

\end{enumerate}
\end{scenariobox}

\vspace{0.10in}

\boxguard[24]
\begin{scenariobox}[A First Look Ahead]{goldacc}
\small
\begin{enumerate}[leftmargin=1.6em, itemsep=6pt, topsep=2pt, start=10]

\item A delivery service charges a flat \$5 for the first $3$ miles and \$2 per mile after that. The
graph shows $F(m)$, the fee in dollars for a delivery of $m$ miles. Answer in plain English where you
do not have the vocabulary yet --- Lesson 2.6 will give you the word.

\begin{center}
\begin{tikzpicture}
\begin{axis}[
    width=9.6cm, height=4.2cm,
    xlabel={$m$ (miles)}, ylabel={$F$ (fee, \$)},
    xmin=0, xmax=8, ymin=0, ymax=16,
    xtick={0,1,2,3,4,5,6,7,8}, ytick={0,5,10,15}, minor y tick num=4,
    grid=both, grid style={linegray!45}, minor grid style={linegray!30},
    axis lines=left,
    tick label style={font=\tiny}, label style={font=\scriptsize}]
  \addplot[royal, very thick] coordinates {(0,5) (3,5) (8,15)};
  \addplot[keyred, only marks, mark=*, mark size=1.7pt] coordinates {(2,5) (3,5) (6,11)};
\end{axis}
\end{tikzpicture}
\end{center}

\begin{enumerate}[label=(\alph*), leftmargin=1.6em, itemsep=5pt, topsep=4pt]
  \item Read $F(2)$ and $F(6)$ off the graph, and say what each one means.

\ansline{$F(2)=5$: a two-mile delivery costs \$5, the flat rate.}\\
\ansline{$F(6)=11$: a six-mile delivery costs \$11 --- \$5 plus \$2 for each of the three extra miles.}\\

  \item This graph takes \emph{two} rules to describe: one for the first stretch and one for the
        rest. Write the two sentences, and say which input is the hand-off point.

\ansline{From $m=0$ to $m=3$ the fee is constant at \$5. After $m=3$ it rises \$2 for every mile.}\\
\ansline{The hand-off is at $m=3$ miles, where the flat stretch ends and the climb begins.}\\

  \item Describe the end behavior as $m$ runs off to the right, and say what it means about the fee.

\ansline{As $m \to \infty$ the fee climbs without bound --- it rises \$2 per mile forever.}\\
\ansline{There is no longest delivery that costs the most; a farther one always costs more.}\\

  \item Does $F$ have a horizontal asymptote? Answer yes or no and give the reason.

\ansline{No. A horizontal asymptote needs the outputs to level off toward a number, and these}\\
\ansline{outputs pass every dollar amount you could name. Climbing without bound rules one out.}\\
\end{enumerate}

\end{enumerate}
\end{scenariobox}

\vspace{0.10in}

\boxguard[26]
\begin{extensionbox}
\small
\textbf{Same end behavior, very different graphs.} Both $y=2^x$ and $y=x^2$ climb without bound as
$x$ runs off to the right. That does not mean they climb the same way.

\begin{enumerate}[label=(\alph*), leftmargin=1.6em, itemsep=5pt, topsep=4pt]
  \item Complete the table.
\smallskip

\renewcommand{\arraystretch}{1.3}
\begin{tabular}{@{}|l|c|c|c|c|c|c|@{}}
\hline
$x$ & $1$ & $2$ & $3$ & $4$ & $5$ & $10$ \\ \hline
$2^x$ & \ans{$2$} & \ans{$4$} & \ans{$8$} & \ans{$16$} & \ans{$32$} &
    \ans{$1024$} \\ \hline
$x^2$ & \ans{$1$} & \ans{$4$} & \ans{$9$} & \ans{$16$} & \ans{$25$} &
    \ans{$100$} \\
\hline
\end{tabular}

  \item Which is bigger at $x=3$? Which is bigger at $x=10$? Name every input in your table where the
        two are equal.

\ansline{At $x=3$: $x^2$ is bigger ($9 > 8$). At $x=10$: $2^x$ is bigger, and not by a little ($1024 > 100$).}\\
\ansline{They are equal at $x=2$ (both $4$) and at $x=4$ (both $16$).}\\

  \item Both functions have the same end behavior on the right. Write one sentence saying what end
        behavior \emph{does} tell you and one saying what it does \emph{not}.

\ansline{End behavior tells you the \emph{direction} the outputs go at the far ends: both climb}\\
\ansline{without bound, so neither has a horizontal asymptote or an absolute maximum.}\\
\ansline{It does \emph{not} tell you how fast, or which one is ahead --- $x^2$ leads at $x=3$ and loses badly by $x=10$.}\\
\end{enumerate}
\end{extensionbox}

\vspace{0.10in}

\begin{remindbox}
\small
\textbf{Next: Lesson 2.6 --- Piecewise-Defined Functions} (\texttt{AFDA.AF.2}). Item 10's delivery fee
needed \emph{two} rules on two different stretches of the domain, with a hand-off at $m=3$. Next
lesson names that a \textbf{piecewise-defined} function and shows you how to read one, write one, and
find its characteristics --- including the ones you learned today.
\end{remindbox}

\end{document}
